\chapter{Building a Complete Application}
\label{ch:complete-app}

This chapter walks through building a real application from scratch, applying everything covered in the book.  The application is a project tracking system that manages employees, projects, and assignments with full CRUD operations, a data grid, popups, and reports.

\section{Requirements}
\label{sec:app-requirements}

The system tracks:

\begin{itemize}
  \item \textbf{Departments} --- organizational units.
  \item \textbf{Employees} --- name, department, active status.
  \item \textbf{Projects} --- name, start date, status (Active, On Hold, Completed).
  \item \textbf{Assignments} --- which employees are assigned to which projects, with role and hours allocated.
\end{itemize}

\noindent
Users can:
\begin{itemize}
  \item List, create, edit, and delete records for each entity.
  \item Search and filter employee and project lists.
  \item View a project's roster of assigned employees.
  \item Generate a PDF project status report.
  \item Export assignment data to CSV.
\end{itemize}

\section{Database Schema}
\label{sec:app-schema}

Place this SQL in a setup script or run it against your database:

\begin{lstlisting}[style=sql]
CREATE TABLE department (
    department_id  CHAR(36) PRIMARY KEY,
    name           VARCHAR(100) NOT NULL
);

CREATE TABLE employee (
    employee_id    CHAR(36) PRIMARY KEY,
    first_name     VARCHAR(50) NOT NULL,
    last_name      VARCHAR(50) NOT NULL,
    department_id  CHAR(36)
        REFERENCES department(department_id),
    hire_date      INTEGER,
    active         CHAR(1) DEFAULT 'Y'
);

CREATE TABLE project (
    project_id     CHAR(36) PRIMARY KEY,
    name           VARCHAR(100) NOT NULL,
    start_date     INTEGER,
    status         VARCHAR(20) DEFAULT 'Active'
);

CREATE TABLE assignment (
    assignment_id  CHAR(36) PRIMARY KEY,
    project_id     CHAR(36)
        REFERENCES project(project_id),
    employee_id    CHAR(36)
        REFERENCES employee(employee_id),
    role           VARCHAR(50),
    hours_allocated NUMERIC(6,1)
);
\end{lstlisting}

\noindent
Primary keys use UUID strings (36 characters) generated by the backend.  Dates are stored as YYYYMMDD integers, consistent with Kiss conventions.

\section{Backend Services}
\label{sec:app-services}

Each entity gets its own Groovy service in \texttt{src/main/backend/services/}.

\subsection{DepartmentService}

\begin{lstlisting}[style=groovy]
// services/DepartmentService.groovy
import org.kissweb.database.*
import org.kissweb.restServer.ProcessServlet
import org.kissweb.json.JSONObject

class DepartmentService {

    void list(JSONObject injson, JSONObject outjson,
              Connection db, ProcessServlet servlet) {
        outjson.put("departments",
            db.fetchAllJSON(
                "select * from department order by name"))
    }

    void newRecord(JSONObject injson, JSONObject outjson,
                   Connection db, ProcessServlet servlet) {
        Record rec = db.newRecord("department")
        rec.set("department_id",
            UUID.randomUUID().toString())
        rec.set("name", injson.getString("name"))
        rec.addRecord()
        db.commit()
    }

    void save(JSONObject injson, JSONObject outjson,
              Connection db, ProcessServlet servlet) {
        Record rec = db.fetchOne(
            "select * from department " +
            "where department_id = ?",
            injson.getString("departmentId"))
        rec.set("name", injson.getString("name"))
        rec.update()
        db.commit()
    }

    void delete(JSONObject injson, JSONObject outjson,
                Connection db, ProcessServlet servlet) {
        db.execute(
            "delete from department " +
            "where department_id = ?",
            injson.getString("departmentId"))
        db.commit()
    }
}
\end{lstlisting}

\subsection{EmployeeService}

\begin{lstlisting}[style=groovy]
// services/EmployeeService.groovy
import org.kissweb.database.*
import org.kissweb.restServer.ProcessServlet
import org.kissweb.json.JSONObject

class EmployeeService {

    void search(JSONObject injson, JSONObject outjson,
                Connection db, ProcessServlet servlet) {
        QueryBuilder qb = db.newQueryBuilder()
        qb.select("employee.employee_id",
            "employee.first_name",
            "employee.last_name", "employee.hire_date",
            "employee.active",
            "department.name as department_name")
        qb.leftJoin("employee", "department_id",
            "department", "department_id")

        // Dynamic search filters
        String lastName = injson.optString("lastName", "")
        if (!lastName.isEmpty())
            qb.where("employee.last_name like ?",
                lastName + "%")

        String active = injson.optString("active", "")
        if (!active.isEmpty())
            qb.where("employee.active = ?", active)

        qb.orderBy("employee.last_name")
        qb.orderBy("employee.first_name")

        outjson.put("employees", qb.fetchAllJSON())
    }

    void get(JSONObject injson, JSONObject outjson,
             Connection db, ProcessServlet servlet) {
        Record rec = db.fetchOne(
            "select * from employee " +
            "where employee_id = ?",
            injson.getString("employeeId"))
        outjson.put("employee", rec.toJSON())
    }

    void newRecord(JSONObject injson, JSONObject outjson,
                   Connection db, ProcessServlet servlet) {
        Record rec = db.newRecord("employee")
        rec.set("employee_id",
            UUID.randomUUID().toString())
        rec.set("first_name",
            injson.getString("firstName"))
        rec.set("last_name",
            injson.getString("lastName"))
        rec.set("department_id",
            injson.getString("departmentId"))
        rec.set("hire_date",
            injson.getInt("hireDate"))
        rec.set("active", "Y")
        rec.addRecord()
        db.commit()
    }

    void save(JSONObject injson, JSONObject outjson,
              Connection db, ProcessServlet servlet) {
        Record rec = db.fetchOne(
            "select * from employee " +
            "where employee_id = ?",
            injson.getString("employeeId"))
        rec.set("first_name",
            injson.getString("firstName"))
        rec.set("last_name",
            injson.getString("lastName"))
        rec.set("department_id",
            injson.getString("departmentId"))
        rec.set("hire_date",
            injson.getInt("hireDate"))
        rec.set("active",
            injson.getString("active"))
        rec.update()
        db.commit()
    }

    void delete(JSONObject injson, JSONObject outjson,
                Connection db, ProcessServlet servlet) {
        db.execute(
            "delete from employee where employee_id = ?",
            injson.getString("employeeId"))
        db.commit()
    }
}
\end{lstlisting}

\subsection{ProjectService}

\begin{lstlisting}[style=groovy]
// services/ProjectService.groovy
import org.kissweb.database.*
import org.kissweb.restServer.ProcessServlet
import org.kissweb.json.JSONObject

class ProjectService {

    void list(JSONObject injson, JSONObject outjson,
              Connection db, ProcessServlet servlet) {
        String status = injson.optString("status", "")
        String sql = "select * from project"
        if (!status.isEmpty())
            sql += " where status = ?"
        sql += " order by name"

        if (!status.isEmpty())
            outjson.put("projects",
                db.fetchAllJSON(sql, status))
        else
            outjson.put("projects",
                db.fetchAllJSON(sql))
    }

    void get(JSONObject injson, JSONObject outjson,
             Connection db, ProcessServlet servlet) {
        Record rec = db.fetchOne(
            "select * from project " +
            "where project_id = ?",
            injson.getString("projectId"))
        outjson.put("project", rec.toJSON())
    }

    void newRecord(JSONObject injson, JSONObject outjson,
                   Connection db, ProcessServlet servlet) {
        Record rec = db.newRecord("project")
        rec.set("project_id",
            UUID.randomUUID().toString())
        rec.set("name", injson.getString("name"))
        rec.set("start_date",
            injson.getInt("startDate"))
        rec.set("status",
            injson.optString("status", "Active"))
        rec.addRecord()
        db.commit()
    }

    void save(JSONObject injson, JSONObject outjson,
              Connection db, ProcessServlet servlet) {
        Record rec = db.fetchOne(
            "select * from project " +
            "where project_id = ?",
            injson.getString("projectId"))
        rec.set("name", injson.getString("name"))
        rec.set("start_date",
            injson.getInt("startDate"))
        rec.set("status", injson.getString("status"))
        rec.update()
        db.commit()
    }

    void delete(JSONObject injson, JSONObject outjson,
                Connection db, ProcessServlet servlet) {
        // Delete assignments first (referential integrity)
        db.execute(
            "delete from assignment " +
            "where project_id = ?",
            injson.getString("projectId"))
        db.execute(
            "delete from project " +
            "where project_id = ?",
            injson.getString("projectId"))
        db.commit()
    }
}
\end{lstlisting}

\subsection{AssignmentService}

\begin{lstlisting}[style=groovy]
// services/AssignmentService.groovy
import org.kissweb.database.*
import org.kissweb.restServer.ProcessServlet
import org.kissweb.json.JSONObject

class AssignmentService {

    void listForProject(JSONObject injson,
                        JSONObject outjson,
                        Connection db,
                        ProcessServlet servlet) {
        QueryBuilder qb = db.newQueryBuilder()
        qb.select("assignment.assignment_id",
            "assignment.role",
            "assignment.hours_allocated",
            "employee.first_name", "employee.last_name")
        qb.leftJoin("assignment", "employee_id",
            "employee", "employee_id")
        qb.where("assignment.project_id = ?",
            injson.getString("projectId"))
        qb.orderBy("employee.last_name")

        outjson.put("assignments", qb.fetchAllJSON())
    }

    void newRecord(JSONObject injson, JSONObject outjson,
                   Connection db, ProcessServlet servlet) {
        Record rec = db.newRecord("assignment")
        rec.set("assignment_id",
            UUID.randomUUID().toString())
        rec.set("project_id",
            injson.getString("projectId"))
        rec.set("employee_id",
            injson.getString("employeeId"))
        rec.set("role", injson.getString("role"))
        rec.set("hours_allocated",
            injson.getDouble("hoursAllocated"))
        rec.addRecord()
        db.commit()
    }

    void save(JSONObject injson, JSONObject outjson,
              Connection db, ProcessServlet servlet) {
        Record rec = db.fetchOne(
            "select * from assignment " +
            "where assignment_id = ?",
            injson.getString("assignmentId"))
        rec.set("role", injson.getString("role"))
        rec.set("hours_allocated",
            injson.getDouble("hoursAllocated"))
        rec.update()
        db.commit()
    }

    void delete(JSONObject injson, JSONObject outjson,
                Connection db, ProcessServlet servlet) {
        db.execute(
            "delete from assignment " +
            "where assignment_id = ?",
            injson.getString("assignmentId"))
        db.commit()
    }
}
\end{lstlisting}

\section{Frontend Screens}
\label{sec:app-screens}

The frontend consists of HTML screens and their companion JavaScript files in \texttt{src/main/frontend/screens/}.

\subsection{Employee List Screen}

The employee list demonstrates the standard Kiss pattern: a search form, a data grid, and an edit popup.

\begin{lstlisting}[style=html]
<!-- screens/EmployeeList.html -->
<div>
    <h2>Employees</h2>

    <div style="margin-bottom: 10px;">
        <text-input id="searchLastName"
            style="width:200px;"></text-input>
        <drop-down id="searchActive"
            style="width:120px;"></drop-down>
        <push-button id="searchBtn">Search</push-button>
        <push-button id="newBtn">New Employee</push-button>
    </div>

    <div id="employeeGrid"
         style="width:100%; height:400px;"></div>
</div>

<popup id="employeePopup" width="500px" height="380px">
    <popup-title>Employee Detail</popup-title>
    <popup-body>
        <div style="display:grid;
                    grid-template-columns: 120px 1fr;
                    gap: 8px;">
            <label>First Name:</label>
            <text-input id="epFirstName"
                style="width:200px;"></text-input>

            <label>Last Name:</label>
            <text-input id="epLastName"
                style="width:200px;"></text-input>

            <label>Department:</label>
            <drop-down id="epDepartment"
                style="width:200px;"></drop-down>

            <label>Hire Date:</label>
            <date-input id="epHireDate"
                style="width:150px;"></date-input>

            <label>Active:</label>
            <drop-down id="epActive"
                style="width:100px;"></drop-down>
        </div>
        <div style="margin-top:15px; text-align:right;">
            <push-button id="epSave">Save</push-button>
            <push-button id="epCancel">Cancel</push-button>
        </div>
    </popup-body>
</popup>
\end{lstlisting}

\begin{lstlisting}[style=javascript]
// screens/EmployeeList.js
const WS = 'services.EmployeeService';
let grid;
let editMode;  // 'new' or 'edit'
let editId;    // employee_id being edited

// Initialize screen
$$('searchActive').add('', '(all)');
$$('searchActive').add('Y', 'Active');
$$('searchActive').add('N', 'Inactive');
$$('searchActive').setValue('Y');

// Set up the grid
grid = new AGGrid('employeeGrid', [
    { headerName: 'Last Name',
      field: 'last_name', width: 150 },
    { headerName: 'First Name',
      field: 'first_name', width: 150 },
    { headerName: 'Department',
      field: 'department_name', width: 150 },
    { headerName: 'Hire Date',
      field: 'hire_date', width: 120,
      valueFormatter: AGGrid.date },
    { headerName: 'Active',
      field: 'active', width: 80 }
], 'employeeGrid');
grid.show();

// Search
$$('searchBtn').onclick(search);
search();  // initial load

async function search() {
    const res = await Server.call(WS, 'search', {
        lastName:
            $$('searchLastName').getValue(),
        active: $$('searchActive').getValue()
    });
    if (res._Success)
        grid.addRecords(res.employees);
}

// Grid double-click to edit
grid.setOnRowDoubleClicked(openEdit);

async function openEdit() {
    const row = grid.getSelectedRow();
    if (!row) return;

    editMode = 'edit';
    editId = row.employee_id;

    const res = await Server.call(WS, 'get',
        { employeeId: editId });
    if (!res._Success) return;

    const emp = res.employee;
    await loadDepartments();
    $$('epFirstName').setValue(emp.first_name);
    $$('epLastName').setValue(emp.last_name);
    $$('epDepartment').setValue(emp.department_id);
    $$('epHireDate').setValue(emp.hire_date);

    $$('epActive').clear();
    $$('epActive').add('Y', 'Yes');
    $$('epActive').add('N', 'No');
    $$('epActive').setValue(emp.active);

    Utils.popup_open('employeePopup');
}

// New employee
$$('newBtn').onclick(async () => {
    editMode = 'new';
    editId = null;

    await loadDepartments();
    $$('epFirstName').clear();
    $$('epLastName').clear();
    $$('epHireDate').clear();

    $$('epActive').clear();
    $$('epActive').add('Y', 'Yes');
    $$('epActive').add('N', 'No');
    $$('epActive').setValue('Y');

    Utils.popup_open('employeePopup');
});

// Save
$$('epSave').onclick(async () => {
    if ($$('epFirstName').isError('First name'))
        return;
    if ($$('epLastName').isError('Last name'))
        return;

    const data = {
        firstName: $$('epFirstName').getValue(),
        lastName: $$('epLastName').getValue(),
        departmentId: $$('epDepartment').getValue(),
        hireDate: $$('epHireDate').getValue(),
        active: $$('epActive').getValue()
    };

    let res;
    if (editMode === 'new') {
        res = await Server.call(
            WS, 'newRecord', data);
    } else {
        data.employeeId = editId;
        res = await Server.call(WS, 'save', data);
    }

    if (res._Success) {
        Utils.popup_close('employeePopup');
        search();
    }
});

// Cancel
$$('epCancel').onclick(() => {
    Utils.popup_close('employeePopup');
});

// Helper: load department dropdown
async function loadDepartments() {
    const res = await Server.call(
        'services.DepartmentService', 'list');
    if (!res._Success) return;

    $$('epDepartment').clear();
    $$('epDepartment').add('', '(select)');
    for (const dept of res.departments) {
        $$('epDepartment').add(
            dept.department_id, dept.name);
    }
}
\end{lstlisting}

\noindent
This screen demonstrates several key Kiss patterns:

\begin{itemize}
  \item Dynamic search with \texttt{QueryBuilder} on the backend.
  \item AGGrid for data display with built-in date formatting.
  \item Popup reuse for both new and edit operations.
  \item Drop-down population from a related service.
  \item Validation with \texttt{isError()} before saving.
\end{itemize}

\subsection{Project Dashboard}

The project dashboard shows the project list and, when a project is selected, its assigned employees in a detail grid:

\begin{lstlisting}[style=html]
<!-- screens/ProjectDashboard.html -->
<div>
    <h2>Projects</h2>
    <div style="margin-bottom: 10px;">
        <drop-down id="statusFilter"
            style="width:150px;"></drop-down>
        <push-button id="filterBtn">Filter</push-button>
        <push-button id="newProjectBtn">
            New Project</push-button>
        <push-button id="reportBtn">
            Status Report</push-button>
        <push-button id="exportBtn">
            Export CSV</push-button>
    </div>
    <div id="projectGrid"
         style="width:100%; height:250px;"></div>

    <h3 id="rosterTitle"
        style="margin-top:15px;">Assignments</h3>
    <div id="assignmentGrid"
         style="width:100%; height:200px;"></div>
</div>
\end{lstlisting}

\begin{lstlisting}[style=javascript]
// screens/ProjectDashboard.js (abbreviated)
const ProjWS = 'services.ProjectService';
const AsgnWS = 'services.AssignmentService';
let projectGrid, assignmentGrid;
let selectedProjectId = null;

// Status filter
$$('statusFilter').add('', '(all)');
$$('statusFilter').add('Active', 'Active');
$$('statusFilter').add('On Hold', 'On Hold');
$$('statusFilter').add('Completed', 'Completed');
$$('statusFilter').setValue('Active');

// Project grid
projectGrid = new AGGrid('projectGrid', [
    { headerName: 'Project', field: 'name',
      width: 200 },
    { headerName: 'Start Date', field: 'start_date',
      width: 120, valueFormatter: AGGrid.date },
    { headerName: 'Status', field: 'status',
      width: 120 }
], 'projectGrid');
projectGrid.show();

// Assignment grid
assignmentGrid = new AGGrid('assignmentGrid', [
    { headerName: 'Employee',
      valueGetter: (p) =>
          p.data.last_name + ', ' + p.data.first_name,
      width: 200 },
    { headerName: 'Role', field: 'role', width: 150 },
    { headerName: 'Hours', field: 'hours_allocated',
      width: 100, type: 'numericColumn' }
], 'assignmentGrid');
assignmentGrid.show();

// Load projects
$$('filterBtn').onclick(loadProjects);
loadProjects();

async function loadProjects() {
    const res = await Server.call(ProjWS, 'list',
        { status: $$('statusFilter').getValue() });
    if (res._Success)
        projectGrid.addRecords(res.projects);
}

// When a project is selected, load its assignments
projectGrid.setOnSelectionChanged(() => {
    const row = projectGrid.getSelectedRow();
    if (row) {
        selectedProjectId = row.project_id;
        loadAssignments(selectedProjectId);
    }
});

async function loadAssignments(projectId) {
    const res = await Server.call(
        AsgnWS, 'listForProject',
        { projectId: projectId });
    if (res._Success)
        assignmentGrid.addRecords(res.assignments);
}
\end{lstlisting}

\section{Reports}
\label{sec:app-reports}

\subsection{PDF Status Report}

Add a report method to \texttt{ProjectService}:

\begin{lstlisting}[style=groovy]
void statusReport(JSONObject injson, JSONObject outjson,
                  Connection db,
                  ProcessServlet servlet) {
    Groff rpt = new Groff("ProjectStatus-",
        "Project Status Report", false)

    rpt.startTable("Lw(2.5i) Lw(1i) Lw(1i) Nw(1i)")
    rpt.column("Project")
    rpt.column("Start Date")
    rpt.column("Status")
    rpt.column("Hours")
    rpt.endTitle()

    List<Record> projects = db.fetchAll("""
        select p.name, p.start_date, p.status,
            coalesce(sum(a.hours_allocated), 0)
                as total_hours
        from project p
        left join assignment a
            on p.project_id = a.project_id
        group by p.name, p.start_date, p.status
        order by p.name
    """)

    for (Record r : projects) {
        rpt.column(r.getString("name"))
        rpt.dateColumn(r.getInt("start_date"))
        rpt.column(r.getString("status"))
        rpt.column(r.getDouble("total_hours"), "C", 1)
    }

    outjson.put("reportUrl", rpt.process())
}
\end{lstlisting}

\noindent
On the frontend:

\begin{lstlisting}[style=javascript]
$$('reportBtn').onclick(async () => {
    const res = await Server.call(
        ProjWS, 'statusReport');
    if (res._Success)
        Utils.showReport(res.reportUrl);
});
\end{lstlisting}

\subsection{CSV Export}

\begin{lstlisting}[style=groovy]
void exportAssignments(JSONObject injson,
                       JSONObject outjson,
                       Connection db,
                       ProcessServlet servlet) {
    File f = FileUtils.createReportFile(
        "Assignments-", ".csv")
    DelimitedFileWriter dfw =
        new DelimitedFileWriter(f.getAbsolutePath())

    // Header
    dfw.writeField("Project")
    dfw.writeField("Employee")
    dfw.writeField("Role")
    dfw.writeField("Hours")
    dfw.endRecord()

    List<Record> recs = db.fetchAll("""
        select p.name as project_name,
            e.last_name || ', ' || e.first_name
                as employee_name,
            a.role, a.hours_allocated
        from assignment a
        join project p
            on a.project_id = p.project_id
        join employee e
            on a.employee_id = e.employee_id
        order by p.name, e.last_name
    """)

    for (Record r : recs) {
        dfw.writeField(r.getString("project_name"))
        dfw.writeField(r.getString("employee_name"))
        dfw.writeField(r.getString("role"))
        dfw.writeField(r.getDouble("hours_allocated"))
        dfw.endRecord()
    }

    dfw.close()
    outjson.put("exportUrl", FileUtils.getHTTPPath(f))
}
\end{lstlisting}

\noindent
On the frontend:

\begin{lstlisting}[style=javascript]
$$('exportBtn').onclick(async () => {
    const res = await Server.call(
        ProjWS, 'exportAssignments');
    if (res._Success)
        Utils.showReport(res.exportUrl);
});
\end{lstlisting}

\section{Navigation}
\label{sec:app-navigation}

After login, set up a simple navigation framework in \texttt{index.js}:

\begin{lstlisting}[style=javascript]
// After successful login
function showMainScreen() {
    Utils.loadPage('screens/MainLayout', 'appBody');
}
\end{lstlisting}

\begin{lstlisting}[style=html]
<!-- screens/MainLayout.html -->
<div style="display:flex; gap:10px;
            padding:10px; background:#f0f0f0;">
    <push-button id="navEmployees">
        Employees</push-button>
    <push-button id="navProjects">
        Projects</push-button>
</div>
<div id="contentArea"></div>
\end{lstlisting}

\begin{lstlisting}[style=javascript]
// screens/MainLayout.js
$$('navEmployees').onclick(() => {
    Utils.loadPage('screens/EmployeeList',
        'contentArea');
});

$$('navProjects').onclick(() => {
    Utils.loadPage('screens/ProjectDashboard',
        'contentArea');
});

// Load default screen
Utils.loadPage('screens/ProjectDashboard',
    'contentArea');
\end{lstlisting}

\section{Lessons Learned}
\label{sec:app-lessons}

Building this application illustrates several key patterns:

\begin{description}[style=nextline]
  \item[Service simplicity] Each service method is typically 10--30 lines.  The framework handles authentication, connection management, and JSON serialization.

  \item[Popup reuse] A single popup serves both ``new'' and ``edit'' modes, controlled by a flag variable.  This halves the amount of UI code.

  \item[QueryBuilder for search] Dynamic search screens benefit from \texttt{QueryBuilder}'s conditional \texttt{where()} clauses.  Each filter adds a condition only when the user provides a value.

  \item[Grid as the center] Most list screens follow the same pattern: search form at the top, grid in the middle, edit popup on double-click, action buttons that call backend methods.

  \item[Reports are services] PDF reports and CSV exports are just service methods that return a URL.  The frontend opens the URL in a new tab.

  \item[UUID primary keys] Generated on the backend with \texttt{UUID.randomUUID().toString()}, avoiding database-specific auto-increment mechanisms.

  \item[Commit explicitly] Every service method that modifies data ends with \texttt{db.commit()}.  The framework does not auto-commit.
\end{description}

\begin{tipbox}
Start with the database schema, then build the backend services, then the frontend screens.  This bottom-up approach ensures each layer has a solid foundation before the next layer is built.
\end{tipbox}
